\documentclass[../../../include/open-logic-section]{subfiles}

\begin{document}

\olfileid[tr]{sth}{story}{urelements}
\olsection{Urelementler Olsun mu?}

Sonraki birkaç bölümde birikimli-yinelemeli küme anlayışından aksiyomlar
çıkarmaya çalışacağız. Fakat ilerlemeden önce \emph{urelementler} hakkında
biraz daha konuşmalıyız.

\olref[approach]{sec}'deki resim, yalnızca eski \emph{kümelerden} yeni
kümeler oluşturmamıza izin verdi. Oysa belirli \emph{küme olmayanların}---
inekler, domuzlar, kum taneleri ya da başka şeylerin---kümelerin !!{element}s{}ı
olmasına izin vermek isteyebiliriz. Bu durumda belirli temel elemanlarla,
\emph{urelementlerle} başlar; her $S$ evresinde urelementlerden ya da
$S$'den önceki evrelerde oluşturulmuş kümelerden (istenen birleşimle) meydana
gelen ``bütün olası'' kümeleri oluşturduğumuzu söyleriz. Ortaya çıkan resim
daha çok şöyle görünür:
\begin{center}
  \begin{tikzpicture}[scale=0.6]
  \tikzset{cut_here/.style={densely dotted}}
  \draw (-4,6) -- (-1, 0) -- (1,0) -- (4,6);
  \draw[cut_here] (-5,8)--(-4,6);
  \draw[cut_here] (5,8)--(4,6);
  \draw(-1.5, 1)--(1.5, 1);
  \draw(-2, 2)--(2, 2);
  \draw(-2.5, 3)--(2.5,3);
  \draw(-3, 4)--(3, 4);
  \draw(-3.5, 5)--(3.5, 5);
  \draw(-4, 6)--(4, 6);
  \node[label] at (2, 0) {\small 0};
  \node[label] at (2.5, 1) {\small 1};
  \node[label] at (3, 2) {\small 2};
  \node[label] at (3.5, 3) {\small 3};
  \node[label] at (4, 4) {\small 4};
  \node[label] at (4.5, 5) {\small 5};
  \node[label] at (5, 6) {\small 6};
  \end{tikzpicture}
\end{center}
Şimdi bir karar vermeliyiz: \emph{Urelementlere izin verelim mi?}

Felsefi açıdan urelementleri kuramımıza katmak anlamlıdır. Başlıca neden,
küme teorimizi \emph{uygulanabilir} kılmaktır. Bunu açıklamak için
\olref[sfr][siz][]{chap}'de iki $A$ ve~$B$ kümesinin, ancak ve ancak
aralarında bir bire bir örten fonksiyon varsa aynı büyüklükte olduğunu,
yani $\cardeq{A}{B}$ olduğunu söylediğimizi hatırlayın. Tarladaki ineklerle
ağıldaki domuzlar birer küme oluşturuyorsa ``inek sayısı domuz sayısına
eşittir'' savını küme teorisiyle ele alabiliriz. Fakat urelementleri
yasaklarsak ineklerle domuzlar küme \emph{oluşturmaz} ve bu küme-kuramsal
çözümleme kullanılamaz. Hatta küme teorisini kümelerin kendisi dışındaki
herhangi bir şeye doğrudan uygulama olanağımız kalmaz. (Urelementleri dâhil
etmenin başka nedenleri için bkz. \citealt[pp.~vi, 24, 50--1]{Potter2004}.)

Matematikte ise urelementlere izin vermek oldukça enderdir. Bunun bir nedeni,
küme teorisini urelementler olmadan ifade etmenin \emph{çok az} daha kolay
olmasıdır. Fakat bazen urelementleri küme teorisinden dışlamak için daha ilginç
gerekçeler bulunur:
\begin{quote}
  Küme teorisinin matematiğin temeli olduğu inancıyla uyumlu olarak, yalnızca
  kümelerden söz ederek bütün matematiği yakalayabilmeliyiz; dolayısıyla
  değişkenimiz inekler ve domuzlar gibi nesneler üzerinde dolaşmamalıdır.
  %But if $C$ is a cow, $\{C\}$ is a set, but not a legitimate mathematical object.
  \citep[p.~8]{Kunen1980}
\end{quote}
Öyleyse uygulanabilirliğe odaklanmak urelementleri \emph{dâhil etmeyi};
indirgemeci bir temel amaca (matematiği saf küme teorisine indirgemeye)
odaklanmak ise onları \emph{dışlamayı} düşündürebilir. Biraz tembellik de
urelementleri dışlama yönündedir.

En tembel yolu izleyeceğiz. Yine de bunun kısmen pedagojik bir gerekçesi
vardır. Amacımız sizi küme teorisi felsefesine başlamanız için gereken küme
teorisi öğeleriyle tanıştırmaktır. Bu felsefi yazının büyük bölümü ise
urelementler \emph{olmadan} formüle edilmiş küme teorilerini tartışır. Bu
kitap o yazına oldukça yakın durursa belki daha yararlı olur.

\end{document}
